\section{Signed ray capture and exact canonical arm ceilings}
\label{sec:affine-signed-ray-canonical-caps}

We close the signed-direction and constant-arm losses left by
Sections~\ref{sec:affine-collinear-period-energy} and
\ref{sec:affine-catalogue-weight-overlap}.  Retain the three affine arms
\[
 U(h,k)=1+Rh,\qquad
 V(h,k)=1+R(h+Ck),\qquad
 W(h,k)=1+R(h+Bk),\qquad 0<B<C.
\]
A line through at least two lattice points has a primitive signed direction
$(s,t)$ and a parametrization $(h_0+js,k_0+jt)$.  Put
\begin{equation}\label{eq:affine-signed-coefficients}
 A_U=s,\qquad A_V=s+Ct,\qquad A_W=s+Bt,\qquad
 L=\max(|s|,|t|).
\end{equation}
If the index interval has length $H$ inside a square of side difference $N$,
then $HL\le N$.

For a positive pairwise-coprime divisor label
$\lambda=(d_U,d_V,d_W)$, set
\begin{equation}\label{eq:affine-signed-period-capture}
 r_Z=\frac{d_Z}{\gcd(d_Z,|A_Z|)},\qquad
 T_\lambda=r_Ur_Vr_W,\qquad
 C_\lambda=\prod_Z\gcd(d_Z,|A_Z|),\qquad
 D_\lambda=d_Ud_Vd_W.
\end{equation}

\begin{theorem}[Signed period and strict capture squeeze]
\label{thm:affine-signed-capture-squeeze}
Suppose $\lambda$ occurs $n_\lambda$ times on the segment, and put
$a_\lambda=n_\lambda-1$.  Then
\begin{equation}\label{eq:affine-period-capacity}
 T_\lambda\mid |i-j|\quad\hbox{for two occurrence indices},\qquad
 T_\lambda C_\lambda=D_\lambda,\qquad
 a_\lambda T_\lambda L\le N.
\end{equation}
If $N>0$ and $D_\lambda>N^2$, then
\begin{align}
 a_\lambda NL&<C_\lambda,
 \label{eq:affine-linear-capture-squeeze}\\
 a_\lambda^2T_\lambda L^2&<C_\lambda,\qquad
 a_\lambda^3T_\lambda^2L^3<C_\lambda N,
 \label{eq:affine-higher-capture-squeeze}\\
 (a_\lambda T_\lambda L)^3&<D_\lambda N.
 \label{eq:affine-physical-period-cube}
\end{align}
\end{theorem}

\begin{proof}
Subtract the two arm divisibilities.  Actual occurrence makes every $d_Z$
coprime to $R$, so cancellation gives
$d_Z\mid(i-j)A_Z$.  Dividing by the gcd in
\eqref{eq:affine-signed-period-capture} gives $r_Z\mid|i-j|$.
The $r_Z$ remain pairwise coprime, hence their product divides the gap; the
factorization of $D_\lambda$ is immediate.  All occurrence indices lie in
one residue class modulo $T_\lambda$, so
$a_\lambda T_\lambda\le H$ and then
$a_\lambda T_\lambda L\le N$.

Now
\[
 a_\lambda T_\lambda LN\le N^2
       <T_\lambda C_\lambda.
\]
Cancel $T_\lambda$ to obtain
\eqref{eq:affine-linear-capture-squeeze}.  Squaring the period capacity and
using the same strict inequality gives
$a_\lambda^2T_\lambda L^2<C_\lambda$; multiplication by the capacity proves
the cubic inequality.  Multiplying once more by $T_\lambda$ gives
\eqref{eq:affine-physical-period-cube}.
\end{proof}

This theorem includes all signs and all zero coefficients.  For a non-arm
direction none of $A_U,A_V,A_W$ vanishes.  With
\begin{equation}\label{eq:affine-signed-direction-product}
 P(s,t)=|s|\,|s+Ct|\,|s+Bt|,\qquad K=(B+1)(C+1),
\end{equation}
one has
\begin{equation}\label{eq:affine-signed-product-bound}
 C_\lambda\le P(s,t)\le KL^3.
\end{equation}
The constant is sharp at $s=t=1$.  Consequently every repeated large
non-arm label obeys
\begin{equation}\label{eq:affine-nonarm-caps}
 a_\lambda NL<P(s,t),\qquad
 a_\lambda^3T_\lambda^2<KN,\qquad KL^2>N.
\end{equation}
The last condition follows because $KL^2\le N$ would contradict the first
inequality as soon as $a_\lambda\ge1$.  Thus non-arm repetitions are confined
to large primitive directions and retain an inverse-square period saving.

The canonical relation is $C=B+1$.  There are then three zero-coefficient
primitive directions, up to sign.

\begin{theorem}[Exact canonical arm capture]
\label{thm:affine-exact-three-arm-capture}
For the constant $U,V,W$ directions, respectively, the absolute coefficient
triples, scales, and captures are
\[
\begin{array}{c|c|c|c}
\text{constant arm}&(|A_U|,|A_V|,|A_W|)&L&C_\lambda\\ \hline
U&(0,C,B)&1&d_U\\
V&(C,0,1)&C&d_V\\
W&(B,1,0)&B&d_W.
\end{array}
\]
Hence a large occurring label satisfies the three respective linear caps
\begin{equation}\label{eq:affine-linear-three-arm-caps}
 a_\lambda N<d_U,\qquad
 a_\lambda NC<d_V,\qquad
 a_\lambda NB<d_W.
\end{equation}
\end{theorem}

\begin{proof}
Each component of an actual label divides an arm value congruent to one
modulo $R=\rad(BC)$, and is therefore coprime to $B$ and $C$.
For the $U$ direction,
\[
 C_\lambda
 =\gcd(d_U,0)\gcd(d_V,C)\gcd(d_W,B)
 =d_U.
\]
For the $V$ direction the coefficients are $(C,0,C-B)=(C,0,1)$, giving
\(C_\lambda=d_V\).  The $W$ direction is $(B,-1,0)$ and gives
\(C_\lambda=d_W\).  Substitution into
\eqref{eq:affine-linear-capture-squeeze} proves
\eqref{eq:affine-linear-three-arm-caps}.
\end{proof}

With the preceding canonical point bounds
$3d_U\le C^6$ and $3d_V,3d_W\le C^7$, this gives
\[
 3a_\lambda N<C^6,\qquad
 3a_\lambda NC<C^7,\qquad
 3a_\lambda NB<C^7.
\]
These are linear in the relevant label component and improve the former
square-root treatment of the vertical direction.

We next pass from a shifted occupancy to the full cubic energy.  For every
$n\in\N$, with truncated $a=n-1$,
\begin{equation}\label{eq:affine-optimal-shifted-cube}
                         n^3\le1+7a^3.
\end{equation}
For $n\ge2$, the difference of the right side and $(a+1)^3$ is
$3a(a-1)(2a+1)\ge0$; the cases $n=0,1$ are immediate.  Equality at $n=2$
shows that seven is optimal.

Let $\mathcal K$ be the realized selected canonical kernel triples
$\kappa=(k_U,k_V,k_W)$, with $D_\kappa=k_Uk_Vk_W$.  Define
$\mathcal L$ as the deduplicated union of their downward divisor catalogues,
then retain only labels whose coordinate product exceeds $N^2$.  Let
$n_\lambda$ be the number of selected canonical points realizing the label
$\lambda$.  For
\[
 w_\lambda=\varphi(d_U)\varphi(d_V)\varphi(d_W),\qquad
 W_0=\sum_{\lambda\in\mathcal L}w_\lambda,\qquad
 E=\sum_{\lambda\in\mathcal L}w_\lambda n_\lambda^3,\qquad
 E_{\rm sh}=\sum_{\lambda\in\mathcal L}w_\lambda(n_\lambda-1)^3,
\]
equation \eqref{eq:affine-optimal-shifted-cube} gives
\begin{equation}\label{eq:affine-global-shifted-bridge}
                         E\le W_0+7E_{\rm sh}.
\end{equation}
This charges a singleton once and introduces no number-of-lines factor.

Choose one owner $o(\lambda)\in\mathcal K$ for every deduplicated label.
Euler's divisor identity gives, for each coordinate $Z$,
\begin{equation}\label{eq:affine-owner-third-moment}
 \sum_{\lambda\in\mathcal L}w_\lambda d_Z(\lambda)^3
 \le\sum_{\kappa\in\mathcal K}D_\kappa k_Z(\kappa)^3,
 \qquad
 W_0\le\sum_{\kappa\in\mathcal K}D_\kappa.
\end{equation}
Indeed, labels owned by $\kappa$ form a subset of its full downward
catalogue, whose total totient weight is $D_\kappa$, and
$d_Z^3\le k_Z^3$.

Partition repeated labels into non-arm labels and the three arm directions.
Let
\[
 \mathcal S_{\rm non}
 =\sum_{\substack{\lambda\text{ repeated}\\
                   \lambda\text{ non-arm}}}
   \frac{w_\lambda}{T_\lambda^2}.
\]

\begin{theorem}[Global signed affine energy ceiling]
\label{thm:affine-global-signed-energy-ceiling}
In the canonical setup,
\begin{align}
 E\le{}&\sum_{\kappa\in\mathcal K}D_\kappa
       +7KN\mathcal S_{\rm non}\notag\\
 &+\frac7{N^3}\sum_{\kappa\in\mathcal K}D_\kappa
   \left(k_U^3+\frac{k_V^3}{C^3}+\frac{k_W^3}{B^3}\right).
\label{eq:affine-global-signed-energy-ceiling}
\end{align}
Every term in $\mathcal S_{\rm non}$ is supported on $KL^2>N$.
\end{theorem}

\begin{proof}
The non-arm cubic cap in \eqref{eq:affine-nonarm-caps} gives
$E_{\rm non}\le KN\mathcal S_{\rm non}$.  Cubing
\eqref{eq:affine-linear-three-arm-caps}, multiplying by $w_\lambda$, and
summing gives the three coordinate moments divided by
$N^3,N^3C^3,N^3B^3$.  Enlarge each arm subset to all labels and apply
\eqref{eq:affine-owner-third-moment}.  Add the terms and use
\eqref{eq:affine-global-shifted-bridge}.
\end{proof}

The arm-level contribution is therefore explicit.  The remaining affine
upper-bound task is the non-arm inverse-period sum
$\mathcal S_{\rm non}$, together with the comparison of the singleton
baseline against catalogue-class multiplicities.  These are active gates.

Several complete pressure tests determine the exact boundary.  Replacing
$N^2<D_\lambda$ by $N^2\le D_\lambda$ fails at
$(n,H,T,C_\lambda,D,N,L)=(2,1,1,1,1,1,1)$.  The tuple
\[
 (n,H,T,C_\lambda,D,N,K,L)=(2,3,3,6,18,3,6,1)
\]
satisfies the period, span,
factorization, strict-threshold,
and capture premises, but refutes adding one more factor of $T$ to either
higher capture inequality.  Scaling a primitive arm direction by five gives
full-premise examples showing why primitivity is required, and four separate
primitive examples show why each coefficient-coprimality premise is required.
At occupancy two, coefficient six in
\eqref{eq:affine-optimal-shifted-cube} fails.  These witnesses retire only
the stated strengthenings; actual canonical labels retain the strictness,
primitivity, and coprimality premises.

Catalogue membership is also essential.  In the actual box
$(B,C,M,N,R)=(1,2,10,9,2)$, with $82$ admissible points, the totient weight
of all arm-divisor labels above $N^2$ is $972{,}496$, whereas the owner mass
of the twelve distinct powerful-kernel triples is only $1{,}072$.  Thus
deleting membership in the
selected downward-catalogue union would assert the false inequality
$972{,}496\le1{,}072$.  This full-premise counterexample retires precisely
that enlarged owner statement, while leaving the selected-catalogue theorem
unchanged.

An independent exact-integer audit checks $15{,}840$ signed directions,
$1{,}776{,}807$ cubic ledgers, $2{,}390{,}018$ normalized quadratic ledgers,
$43{,}403$ exact arm-capture cases, and four actual canonical boxes.  Those
boxes contain $4{,}307$ distinct local large labels and $259$ repeated
labels, testing every ray-level identity.  The selected powerful-kernel
subcatalogues for $M=10$ contain $110$ large labels and $12$ repeated labels.
A separate $(B,C,M)=(4,5,30)$ stress box contains $113$ selected large labels,
nine repeated labels, and one repeated non-arm label.  All owner and global
energy checks pass.  These finite audits are pressure tests and supply no
asymptotic inference.

The Lean companion kernel-checks the signed period/capture algebra, all three
canonical arm identities, the sharp constant $K$, the optimal coefficient
seven, the finite weighted owner bounds, the global shifted-energy interface,
and the full-premise boundary countermodels.  It assumes no exceptional-point
lower bound, no catalogue novelty theorem, and no $abc$ statement.
